% Production session of Study 1, in parts for content.tex:
%   Part 1 replaces the \todo{Fill in the production session once run: ...}
%          under \subsection{Design} \label{sec:feas-s1-design}
%   Part 2 replaces the \todo under \subsection{Production Configuration}
%          \label{sec:results-s1-production}
%   Part 3 (optional) is a sentence for \subsection{Summary} \label{sec:results-s1-summary}
% Requires: booktabs, pgfplots (compat=1.18) with the groupplots library, \abs from commands.tex. No siunitx.
% Numbers: python3 scripts/analyse_head_accuracy.py study/head_accuracy_production.csv --moves
%          python3 scripts/analyse_head_accuracy.py study/head_accuracy_20260909.csv --only 45:1.25:none 90:1.25:none 135:1.25:none --trials
%          python3 scripts/analyse_head_accuracy.py study/head_accuracy_production.csv --drop production:11 --trials
%          python3 scripts/analyse_head_accuracy.py study/head_accuracy_20260909.csv --trials
%          python3 scripts/analyse_head_accuracy.py study/head_accuracy_validation.csv --trials
% Figure:  python3 scripts/analyse_head_accuracy.py study/head_accuracy_20260909.csv study/head_accuracy_production.csv --only 45:1.25:none 135:1.25:none --pgf-until 4.5

% ======================================================================
% Part 1 - method
% ======================================================================

The production session was run on 14 September 2026 between 15:48 and 15:55:
five trials at each of the three bearings, at 1.25\,m in quiet, 15 in all, with
the stimulus started by hand as in the main session.
\todo{Confirm that the room, the stimulus recording and its 60\,dB(A) level were those of the main session.}
The harness started the tracker with the default arguments of
\texttt{start\_doa\_tracking()}, read from the function itself rather than
copied, and recorded the resulting configuration with every trial: no discard
after the voice-activity onset, whose own reading counts as the first sample; up
to five readings 37.5\,ms apart within 0.6\,s, and at least two; 0.6\,s after
each move during which voice is ignored; and no filter on the estimated bearing.
\texttt{DOA\_OFFAXIS\_GAIN} was 1.0 and no \texttt{DOA\_CALIBRATION\_PATH} was
set, so the head followed the uncorrected bearing, with the front reference at
90\textdegree{}. As in the main session, the head was not returned home after a
silence during a trial, whereas the running assistant returns it after 5\,s of
silence.

% ======================================================================
% Part 2 - results
% ======================================================================

The head reacted in all 15 trials. Pooled over the three bearings, its final
heading was a median 5.0\textdegree{} (IQR 2.0--6.0\textdegree{}) from the true
bearing, against 5.0\textdegree{} (IQR 0.0--6.0\textdegree{}) in the same three
cells of the main session, and it came to rest a median 1.525\,s (IQR
1.336--1.540\,s) after the array's first voice report, against 3.428\,s (IQR
3.413--3.670\,s). Table~\ref{tab:s1-production} breaks both down by bearing. The
pooled figures conceal three differences, taken in turn below: the production
head reaches its bearing in two moves rather than one, its final error is less
consistent, and in one trial it turned away from the source.

\begin{table}[htb]
  \centering
  \small
  \caption{Study and production configuration at 1.25\,m in quiet, five trials
  per cell, with no detection failures in either. Errors are commanded heading
  minus true bearing; the first column gives the error after the first move.
  Times are median [IQR] from the array's first voice report; \emph{latency} is
  the measure of section~\ref{sec:results-s1-latency}. The study
  configuration made one move per trial, so its two time columns coincide.}
  \label{tab:s1-production}
  \setlength{\tabcolsep}{4pt}
  \begin{tabular}{@{}llrrrll@{}}
    \toprule
    & & \multicolumn{2}{c}{Median $\abs{e}$ (\textdegree)} & Signed $e$ (\textdegree) & \multicolumn{2}{c}{Time (s)} \\
    \cmidrule(lr){3-4} \cmidrule(l){6-7}
    Bearing & Config. & first move & final & mean $\pm$ SD & first move & latency \\
    \midrule
    45\textdegree{}  & study      &  6.0 & 6.0 & $+6.40 \pm 1.14$   & 3.666 [3.662, 3.670] & 3.666 [3.662, 3.670] \\
                     & production & 17.0 & 5.0 & $+4.20 \pm 6.98$   & 0.493 [0.422, 0.525] & 1.538 [1.506, 1.542] \\
    \addlinespace
    90\textdegree{}  & study      &  0.0 & 0.0 & $0.00 \pm 0.00$    & 3.378 [3.122, 4.402] & 3.378 [3.122, 4.402] \\
                     & production &  2.0 & 2.0 & $-1.20 \pm 1.10$   & 0.367 [0.259, 0.367] & 1.250 [1.250, 1.250] \\
    \addlinespace
    135\textdegree{} & study      &  6.0 & 6.0 & $-5.40 \pm 0.89$   & 3.414 [3.414, 3.427] & 3.414 [3.414, 3.427] \\
                     & production & 19.0 & 6.0 & $-31.80 \pm 57.71$ & 0.533 [0.485, 0.539] & 1.534 [1.525, 1.538] \\
    \bottomrule
  \end{tabular}
\end{table}

\subsubsection{Two Moves Instead of One}

Figure~\ref{fig:s1-production} shows how the head gets there. In the study
configuration it stays at home until the estimate exists and then turns once.
In the production configuration the first move was complete a median 0.493\,s
and 0.533\,s after the first voice report at 45\textdegree{} and
135\textdegree{}, but it left the head a median 17.0\textdegree{} and
19.0\textdegree{} from the source. The readings behind that move are taken as
soon as voice activity begins, and in several trials some of them still pointed
close to straight ahead: in trial~3 at 45\textdegree{} all five lay between
86\textdegree{} and 88\textdegree{}, and the head turned by 3\textdegree{}; in
trial~13 at 135\textdegree{} three of the five read 90\textdegree{}, and the head
did not turn at all. Once the 0.6\,s after the move had passed, a second estimate
turned the head again in four of the five trials at each off-axis bearing. At
90\textdegree{} the head turned by 2\textdegree{} in three trials and not at all
in two. Every trial produced a second estimate while the 13.6\,s stimulus was still
playing; an utterance that ends before the second estimate would leave the head at
its first position.

\begin{figure}[htb]
  \centering
  \definecolor{studyc}{HTML}{9A9A9A}
  \definecolor{prodc}{HTML}{2A78D6}
  \begin{tikzpicture}
    \begin{groupplot}[
      group style={group size=2 by 1, horizontal sep=1.2cm},
      scale only axis, width=0.39\textwidth, height=0.36\textwidth,
      xmin=0, xmax=4.5, xtick={0,1,2,3,4},
      tick label style={font=\footnotesize},
      label style={font=\footnotesize},
      title style={font=\footnotesize},
      grid=major, grid style={line width=0.2pt, draw=black!12},
      axis line style={draw=black!45}, tick style={draw=black!45},
      xlabel={Time from first voice report (s)},
    ]
      \nextgroupplot[title={45\textdegree{}}, ymin=-20, ymax=50,
        ytick={-10,0,10,20,30,40,50}, ylabel={Heading error (\textdegree)}]
      \addplot[black!55, line width=0.6pt, domain=0:4.5, samples=2, forget plot] {0};
      % trial 26
      \addplot[studyc, line width=0.8pt, forget plot]
        coordinates {(0.000,45.00) (3.122,45.00) (3.122,39.29) (3.173,39.29) (3.173,33.57) (3.224,33.57) (3.224,27.86) (3.275,27.86) (3.275,22.14) (3.326,22.14) (3.326,16.43) (3.377,16.43) (3.377,10.71) (3.427,10.71) (3.427,5.00) (4.500,5.00)};
      % trial 27
      \addplot[studyc, line width=0.8pt, forget plot]
        coordinates {(0.000,45.00) (3.378,45.00) (3.378,38.67) (3.436,38.67) (3.436,32.33) (3.493,32.33) (3.493,26.00) (3.551,26.00) (3.551,19.67) (3.608,19.67) (3.608,13.33) (3.666,13.33) (3.666,7.00) (4.500,7.00)};
      % trial 28
      \addplot[studyc, line width=0.8pt, forget plot]
        coordinates {(0.000,45.00) (3.378,45.00) (3.378,38.83) (3.435,38.83) (3.435,32.67) (3.492,32.67) (3.492,26.50) (3.548,26.50) (3.548,20.33) (3.605,20.33) (3.605,14.17) (3.662,14.17) (3.662,8.00) (4.500,8.00)};
      % trial 29
      \addplot[studyc, line width=0.8pt, forget plot]
        coordinates {(0.000,45.00) (3.378,45.00) (3.378,38.50) (3.436,38.50) (3.436,32.00) (3.495,32.00) (3.495,25.50) (3.553,25.50) (3.553,19.00) (3.611,19.00) (3.611,12.50) (3.670,12.50) (3.670,6.00) (4.500,6.00)};
      % trial 30
      \addplot[studyc, line width=0.8pt, forget plot]
        coordinates {(0.000,45.00) (3.378,45.00) (3.378,38.50) (3.437,38.50) (3.437,32.00) (3.495,32.00) (3.495,25.50) (3.553,25.50) (3.553,19.00) (3.612,19.00) (3.612,12.50) (3.670,12.50) (3.670,6.00) (4.500,6.00)};
      % trial 1
      \addplot[prodc, line width=0.8pt, forget plot]
        coordinates {(0.000,45.00) (0.258,45.00) (0.258,38.75) (0.309,38.75) (0.309,32.50) (0.360,32.50) (0.360,26.25) (0.411,26.25) (0.411,20.00) (0.462,20.00) (0.462,13.75) (0.512,13.75) (0.512,7.50) (0.563,7.50) (0.563,1.25) (0.614,1.25) (0.614,-5.00) (4.500,-5.00)};
      % trial 2
      \addplot[prodc, line width=0.8pt, forget plot]
        coordinates {(0.000,45.00) (0.258,45.00) (0.258,41.75) (0.313,41.75) (0.313,38.50) (0.367,38.50) (0.367,35.25) (0.422,35.25) (0.422,32.00) (1.314,32.00) (1.314,26.80) (1.371,26.80) (1.371,21.60) (1.428,21.60) (1.428,16.40) (1.485,16.40) (1.485,11.20) (1.542,11.20) (1.542,6.00) (4.500,6.00)};
      % trial 3
      \addplot[prodc, line width=0.8pt, forget plot]
        coordinates {(0.000,45.00) (0.258,45.00) (0.258,44.00) (0.314,44.00) (0.314,43.00) (0.370,43.00) (0.370,42.00) (1.250,42.00) (1.250,36.40) (1.309,36.40) (1.309,30.80) (1.368,30.80) (1.368,25.20) (1.426,25.20) (1.426,19.60) (1.485,19.60) (1.485,14.00) (4.500,14.00)};
      % trial 4
      \addplot[prodc, line width=0.8pt, forget plot]
        coordinates {(0.000,45.00) (0.258,45.00) (0.258,39.40) (0.317,39.40) (0.317,33.80) (0.375,33.80) (0.375,28.20) (0.434,28.20) (0.434,22.60) (0.493,22.60) (0.493,17.00) (1.378,17.00) (1.378,14.00) (1.432,14.00) (1.432,11.00) (1.485,11.00) (1.485,8.00) (1.538,8.00) (1.538,5.00) (4.500,5.00)};
      % trial 5
      \addplot[prodc, line width=0.8pt, forget plot]
        coordinates {(0.000,45.00) (0.258,45.00) (0.258,39.50) (0.311,39.50) (0.311,34.00) (0.365,34.00) (0.365,28.50) (0.418,28.50) (0.418,23.00) (0.471,23.00) (0.471,17.50) (0.525,17.50) (0.525,12.00) (1.458,12.00) (1.458,9.25) (1.510,9.25) (1.510,6.50) (1.562,6.50) (1.562,3.75) (1.614,3.75) (1.614,1.00) (4.500,1.00)};
      \nextgroupplot[title={135\textdegree{}}, ymin=-50, ymax=20,
        ytick={-50,-40,-30,-20,-10,0,10,20},
        legend style={at={(0.97,0.97)}, anchor=north east, font=\footnotesize,
                      draw=none, fill=white, cells={anchor=west}}]
      \addplot[black!55, line width=0.6pt, domain=0:4.5, samples=2, forget plot] {0};
      % trial 36
      \addplot[studyc, line width=0.8pt]
        coordinates {(0.000,-45.00) (3.122,-45.00) (3.122,-38.50) (3.181,-38.50) (3.181,-32.00) (3.239,-32.00) (3.239,-25.50) (3.297,-25.50) (3.297,-19.00) (3.356,-19.00) (3.356,-12.50) (3.414,-12.50) (3.414,-6.00) (4.500,-6.00)};
      \addlegendentry{study configuration}
      % trial 37
      \addplot[studyc, line width=0.8pt, forget plot]
        coordinates {(0.000,-45.00) (3.378,-45.00) (3.378,-39.14) (3.430,-39.14) (3.430,-33.29) (3.482,-33.29) (3.482,-27.43) (3.534,-27.43) (3.534,-21.57) (3.585,-21.57) (3.585,-15.71) (3.637,-15.71) (3.637,-9.86) (3.689,-9.86) (3.689,-4.00) (4.500,-4.00)};
      % trial 38
      \addplot[studyc, line width=0.8pt, forget plot]
        coordinates {(0.000,-45.00) (3.122,-45.00) (3.122,-38.50) (3.180,-38.50) (3.180,-32.00) (3.239,-32.00) (3.239,-25.50) (3.297,-25.50) (3.297,-19.00) (3.355,-19.00) (3.355,-12.50) (3.413,-12.50) (3.413,-6.00) (4.500,-6.00)};
      % trial 39
      \addplot[studyc, line width=0.8pt, forget plot]
        coordinates {(0.000,-45.00) (3.122,-45.00) (3.122,-39.29) (3.173,-39.29) (3.173,-33.57) (3.224,-33.57) (3.224,-27.86) (3.275,-27.86) (3.275,-22.14) (3.325,-22.14) (3.325,-16.43) (3.376,-16.43) (3.376,-10.71) (3.427,-10.71) (3.427,-5.00) (4.500,-5.00)};
      % trial 40
      \addplot[studyc, line width=0.8pt, forget plot]
        coordinates {(0.000,-45.00) (3.122,-45.00) (3.122,-38.50) (3.181,-38.50) (3.181,-32.00) (3.239,-32.00) (3.239,-25.50) (3.297,-25.50) (3.297,-19.00) (3.355,-19.00) (3.355,-12.50) (3.414,-12.50) (3.414,-6.00) (4.500,-6.00)};
      % trial 11
      \addplot[prodc, line width=0.8pt]
        coordinates {(0.000,-45.00) (0.258,-45.00) (0.258,-72.58) (0.309,-72.58) (0.309,-100.17) (0.360,-100.17) (0.360,-127.75) (0.411,-127.75) (0.411,-135.00) (4.500,-135.00)};
      \addlegendentry{production}
      % trial 12
      \addplot[prodc, line width=0.8pt, forget plot]
        coordinates {(0.000,-45.00) (0.258,-45.00) (0.258,-39.17) (0.313,-39.17) (0.313,-33.33) (0.368,-33.33) (0.368,-27.50) (0.423,-27.50) (0.423,-21.67) (0.478,-21.67) (0.478,-15.83) (0.533,-15.83) (0.533,-10.00) (1.410,-10.00) (1.410,-8.67) (1.468,-8.67) (1.468,-7.33) (1.525,-7.33) (1.525,-6.00) (4.500,-6.00)};
      % trial 13
      \addplot[prodc, line width=0.8pt, forget plot]
        coordinates {(0.000,-45.00) (1.138,-45.00) (1.138,-38.83) (1.195,-38.83) (1.195,-32.67) (1.251,-32.67) (1.251,-26.50) (1.308,-26.50) (1.308,-20.33) (1.365,-20.33) (1.365,-14.17) (1.422,-14.17) (1.422,-8.00) (4.500,-8.00)};
      % trial 14
      \addplot[prodc, line width=0.8pt, forget plot]
        coordinates {(0.000,-45.00) (0.258,-45.00) (0.258,-39.80) (0.315,-39.80) (0.315,-34.60) (0.372,-34.60) (0.372,-29.40) (0.428,-29.40) (0.428,-24.20) (0.485,-24.20) (0.485,-19.00) (1.362,-19.00) (1.362,-15.25) (1.419,-15.25) (1.419,-11.50) (1.477,-11.50) (1.477,-7.75) (1.534,-7.75) (1.534,-4.00) (4.500,-4.00)};
      % trial 15
      \addplot[prodc, line width=0.8pt, forget plot]
        coordinates {(0.000,-45.00) (0.260,-45.00) (0.260,-39.00) (0.316,-39.00) (0.316,-33.00) (0.372,-33.00) (0.372,-27.00) (0.428,-27.00) (0.428,-21.00) (0.483,-21.00) (0.483,-15.00) (0.539,-15.00) (0.539,-9.00) (1.426,-9.00) (1.426,-8.00) (1.482,-8.00) (1.482,-7.00) (1.538,-7.00) (1.538,-6.00) (4.500,-6.00)};
      \node[font=\scriptsize, prodc, anchor=south west] at (axis cs:0.45,-50) {trial~11: to $-135$\textdegree{}};
    \end{groupplot}
  \end{tikzpicture}
  \caption{Heading error (commanded heading minus true bearing) over time at
  1.25\,m in quiet, one line per trial, stepped at every servo update and held
  after the last. The head starts at home, 45\textdegree{} from the source. Grey:
  study configuration; blue: production. Trial~11 leaves the plot at the bottom
  and stays at $-135$\textdegree{}. At 90\textdegree{} (not shown) the head turned
  by 2\textdegree{} at most in either configuration.}
  \label{fig:s1-production}
\end{figure}

\subsubsection{Final Error}

Once at rest, the production head was about as close to the source as the study
configuration left it: a median 5.0\textdegree{} against 6.0\textdegree{} at
45\textdegree{}, 6.0\textdegree{} against 6.0\textdegree{} at 135\textdegree{},
and 2.0\textdegree{} against 0.0\textdegree{} at 90\textdegree{}. It was less
consistent. At 45\textdegree{} the final error ranged from $-5$\textdegree{} to
$+14$\textdegree{}, with a standard deviation of 6.98\textdegree{} against
1.14\textdegree{}. At 135\textdegree{}, leaving out the trial described next, the
signed mean was $-6.00 \pm 1.63$\textdegree{} ($n = 4$), against
$-5.40 \pm 0.89$\textdegree{}. With five trials per cell the comparison is
descriptive and was not tested.

\subsubsection{A Turn Away from the Source}

In trial~11 at 135\textdegree{} the head turned to 0\textdegree{}, the limit of
its sweep on the side away from the source, and stayed there: an error of
$-135$\textdegree{}. The readings are expressed relative to the head, with
90\textdegree{} straight ahead, and run from $-90$\textdegree{} to
270\textdegree{}, both of which lie directly behind it. The first estimate
rested on the readings $-90$, $-90$, 119, 119 and 119\textdegree{}. The circular
median (section~\ref{sec:feas-s1-procedure}) unwraps each reading onto the branch
of the circle nearest the first one, which placed 119\textdegree{} at
$-241$\textdegree{}, and the median of the unwrapped values, $-241$\textdegree{},
lies far outside the sweep; the servo clamped it to 0\textdegree{}. The second
estimate began the same way ($-62$, $-62$, 232, 230 and 231\textdegree{}), came
out at $-128$\textdegree{}, and left the head where it was. The study
configuration waits 1.5\,s before sampling and discards estimates outside the
sweep rather than clamping them; production does neither.

A plain median of the same five readings is 119\textdegree{}, which would have
turned the head towards the source. Where it would then have come to rest cannot
be reconstructed, because every later reading depends on the heading. To check
what the choice of median does elsewhere, every move of the three sessions was
rebuilt from its logged readings. The rebuilt circular median reproduced the
recorded estimate in all 133 moves (83 in the main session, 15 in the validation
session and 35 in this one), and a plain median differed from it only in the two
moves of trial~11. None of the readings that fed no move spanned 180\textdegree{}
or more, the condition for the two medians to disagree. The results of the main
and validation sessions are therefore unchanged under a plain median.

\subsubsection{What the Latency Measures}

The latency measure of section~\ref{sec:results-s1-latency} runs to the last
move the tracker issued before the trial ended, whether or not that move changed
the heading. In the study configuration every trial ended after one move, so the
measure is the time to that move. The production tracker keeps estimating while
speech continues, and the measure then includes estimates that leave the head
where it is. At 90\textdegree{} the recorded 1.250\,s is the time of a second
estimate that did not move the head, and the longest value, 3.826\,s in
trial~10, comes from five estimates none of which moved it. Measured instead to
the last update that changed the commanded angle, the off-axis medians are
1.538\,s (IQR 1.485--1.542\,s) at 45\textdegree{} and 1.525\,s (IQR
1.422--1.534\,s) at 135\textdegree{}, close to the recorded values, since at
those bearings the second estimate was a real move in four of five trials.

\subsubsection{Is the Study Configuration Representative?}

The session answers the question it was run for in two parts. For where the head
comes to rest at 1.25\,m in quiet, the study configuration is representative:
the final errors are of the same size, and at both off-axis bearings the signed
mean points towards the front in both configurations, at 135\textdegree{} once
trial~11 is left out. For how the head gets
there, and how reliably, it is not. The deployed head completes a first move a
median 0.367--0.533\,s after the first voice report and is at rest after
1.525\,s rather than 3.428\,s, but its first move is coarse, its final bearing
is less consistent, and in one of 15 trials the combination of sampling at the
onset and a circular median turned it away from the talker. Fifteen trials show
that this failure occurs; they do not estimate how often. Neither distance nor
noise was varied, and both affected the study configuration strongly
(sections~\ref{sec:results-s1-distance} and~\ref{sec:results-s1-noise}).

% ======================================================================
% Part 3 (optional) - for \subsection{Summary} \label{sec:results-s1-summary},
% after "... 5.2\,s under babble."
% ======================================================================

% With the production defaults, at 1.25\,m in quiet, the head came to rest a
% median 1.5\,s after voice onset with a final error of the same size, but it
% reached that bearing through a coarse first move, and in one of 15 trials it
% turned away from the talker (section~\ref{sec:results-s1-production}).
